\section{Experiment P1F-F: non-Gaussian and NLOS-like UWB stress}
\subsection{Purpose and claim boundary}
P1F-F tests whether the frozen literal AACOPF particle-management mechanism has any inherent robustness advantage over conventional particle-filter resampling when the UWB measurement process departs from the Gaussian likelihood assumed by both filters. The experiment deliberately leaves the estimator-side range model unchanged,
\begin{equation}
    p(z_k\mid\bm x_k)\propto
    \exp\!\left[-\frac{(z_k-\hat d_k)^2}{2\sigma_{\mathrm{UWB}}^2}\right],
    \qquad \sigma_{\mathrm{UWB}}=\SI{0.12}{m}.
\end{equation}
No robust likelihood, innovation gate, NLOS detector, measurement rejection rule, or parameter retuning is introduced. Hence P1F-F is not an evaluation of an explicitly robust ranging estimator. It isolates whether conventional resampling or the frozen ACO copying mechanism is more tolerant to measurement-model mismatch.

The AACOPF parameters remain fixed at the P1F-C values
\begin{equation}
    (\alpha,\beta,c_\lambda)=(0.5,0,2).
\end{equation}
They are not retuned for any corruption condition.

\subsection{Controlled mode-selection setting}
The random-annulus initialization used in P1F-D is intentionally not used for this experiment. P1F-D showed that at the manageable literal-AACOPF budget $N_p=1200$, the clean-data experiment is already dominated by sparse initial support of the correct mode. Such failures would confound a measurement-robustness test.

P1F-F therefore returns to the controlled two-mode setting from P1F-C with $N_p=400$ particles and the informative moving auxiliary. Two initial particle clouds are evaluated:
\begin{itemize}
    \item \textbf{balanced:} $50\%$ of particles represent the correct mode and $50\%$ the competing mode;
    \item \textbf{minority-correct:} $20\%$ represent the correct mode and $80\%$ the competing mode.
\end{itemize}
The competing mode is displaced by $90^\circ$ in initial bearing and yaw. Twenty previously unused seeds, 200--219, are evaluated for every condition and both initial clouds.

A correct-mode lock requires the method-specific correct-mode signal to remain at or above $0.9$ to the end of the run for at least \SI{1}{s}. A wrong-mode lock analogously requires the signal to remain at or below $0.1$. For conventional PF, the signal is the correct-mode posterior mass before resampling. For AACOPF, it is the post-transition correct-mode particle fraction. The experiment also retains strict pose success, position/yaw RMSE, post-corruption recovery, PF resampling diagnostics, and AACOPF ancestry-collapse diagnostics.

\subsection{Ranging corruption models}
All stressed runs share the same nominal Gaussian realization with their clean counterpart for a given seed. Corruption is added only after the first sample so that initialization is identical across conditions. The six measurement conditions are:
\begin{enumerate}
    \item \textbf{clean:} nominal Gaussian noise only;
    \item \textbf{sparse impulses:} two isolated errors of magnitude \SI{1.5}{m} with random sign;
    \item \textbf{moderate mixture:} $8\%$ of samples receive an additional zero-mean gross-error term with standard deviation \SI{0.6}{m};
    \item \textbf{severe mixture:} the same contamination construction with gross-error standard deviation \SI{1.2}{m};
    \item \textbf{moderate NLOS burst:} one contiguous \SI{1.0}{s} interval receives a \SI{0.6}{m} positive bias and \SI{0.08}{m} extra jitter;
    \item \textbf{severe NLOS burst:} one contiguous \SI{2.0}{s} interval receives a \SI{1.2}{m} positive bias and \SI{0.12}{m} extra jitter.
\end{enumerate}
The moderate and severe variants within each stochastic family reuse the underlying corruption pattern where applicable, so the severity comparison is paired rather than generated from unrelated realizations.

\subsection{Balanced-cloud results}
Table~\ref{tab:p1ff_balanced} summarizes the balanced initialization. The clean reference is strong for both algorithms, but the corrupted-range experiments consistently favor conventional PF.

\begin{table}[htbp]
    \centering
    \caption{P1F-F balanced-cloud robustness results over 20 matched seeds. Position RMSE is full-horizon RMSE.}
    \label{tab:p1ff_balanced}
    \begin{tabular}{lrrrrrrr}
        \toprule
        Condition & PF lock & ACO lock & PF wrong & ACO wrong & PF pos. & ACO pos. & ACO collapse\\
        & & & & & [m] & [m] & \\
        \midrule
        Clean                 & 1.00 & 0.95 & 0.00 & 0.00 & 1.715 & 2.251 & 0.05\\
        Sparse impulses       & 1.00 & 0.80 & 0.00 & 0.15 & 1.817 & 2.740 & 0.50\\
        Mixture, moderate     & 0.90 & 0.70 & 0.10 & 0.25 & 2.422 & 2.649 & 0.75\\
        Mixture, severe       & 0.95 & 0.65 & 0.05 & 0.35 & 2.095 & 2.954 & 0.90\\
        NLOS burst, moderate  & 0.75 & 0.60 & 0.25 & 0.35 & 3.031 & 3.819 & 0.50\\
        NLOS burst, severe    & 0.75 & 0.60 & 0.25 & 0.40 & 3.071 & 4.443 & 0.45\\
        \bottomrule
    \end{tabular}
\end{table}

The two-impulse condition is particularly informative. Conventional PF retains a $20/20$ correct-mode lock rate and no wrong locks, whereas AACOPF falls to $16/20$ correct locks and produces $3/20$ wrong locks. Because only two measurements are corrupted, this behavior is consistent with a short-lived likelihood error being amplified by the copying transition.

\subsection{Minority-correct results}
The harder minority-correct initialization gives the same qualitative ordering, as shown in Table~\ref{tab:p1ff_minority}.

\begin{table}[htbp]
    \centering
    \caption{P1F-F minority-correct robustness results over 20 matched seeds.}
    \label{tab:p1ff_minority}
    \begin{tabular}{lrrrrrrr}
        \toprule
        Condition & PF lock & ACO lock & PF wrong & ACO wrong & PF pos. & ACO pos. & ACO collapse\\
        & & & & & [m] & [m] & \\
        \midrule
        Clean                 & 1.00 & 0.95 & 0.00 & 0.00 & 2.745 & 3.508 & 0.15\\
        Sparse impulses       & 0.95 & 0.75 & 0.05 & 0.20 & 2.819 & 4.159 & 0.45\\
        Mixture, moderate     & 0.90 & 0.80 & 0.10 & 0.15 & 3.354 & 4.147 & 0.55\\
        Mixture, severe       & 0.85 & 0.80 & 0.10 & 0.20 & 3.007 & 4.550 & 0.70\\
        NLOS burst, moderate  & 0.80 & 0.70 & 0.20 & 0.25 & 3.894 & 4.945 & 0.30\\
        NLOS burst, severe    & 0.75 & 0.65 & 0.25 & 0.25 & 4.169 & 5.205 & 0.45\\
        \bottomrule
    \end{tabular}
\end{table}

Relative to the clean matched seeds, AACOPF loses more correct-lock runs than PF for every stress family. For the severe mixture the reduction in lock fraction is numerically equal ($0.15$) for both algorithms, but AACOPF starts from a lower clean lock rate, has a higher wrong-lock fraction, and has larger position/yaw RMSE.

\subsection{Pooled robustness comparison}
Pooling the five corrupted conditions, two initial clouds, and 20 seeds gives 200 stressed runs per algorithm. Conventional PF achieves
\begin{equation}
    172/200=86.0\% \quad \text{correct locks},
\end{equation}
with $27/200=13.5\%$ wrong locks. Frozen AACOPF achieves only
\begin{equation}
    141/200=70.5\% \quad \text{correct locks},
\end{equation}
with $51/200=25.5\%$ wrong locks. AACOPF has a lower correct-lock fraction than conventional PF in every one of the ten stressed cells.

Pooling by corruption family gives correct-lock fractions of $97.5\%/77.5\%$ for sparse impulses, $90.0\%/73.75\%$ for mixture contamination, and $76.25\%/63.75\%$ for NLOS bursts, where each pair is conventional PF/AACOPF. The corresponding fractions that recover the correct-mode signal after the final corruption are $97.5\%/82.5\%$, $55.0\%/42.5\%$, and $71.25\%/65.0\%$.

\subsection{Ancestry collapse under corrupted likelihood evidence}
Catastrophic AACOPF ancestry collapse occurs in $111/200=55.5\%$ of stressed runs. Among these runs, $40/111=36.0\%$ finish in wrong-mode lock. Without catastrophic collapse, the corresponding rate is $11/89=12.4\%$. This is a descriptive association and does not establish causality, but it is consistent with the P1F-B/P1F-C mechanism diagnosis: if corrupted measurement evidence temporarily favors a wrong hypothesis, aggressive cloning can rapidly remove alternative support and make recovery difficult.

The result is important because it clarifies that the ACO transition only sees posterior weight differences. It has no mechanism for distinguishing a trustworthy high-weight particle from one made temporarily attractive by an outlier or NLOS-biased range.

\subsection{P1F-F interpretation and decision}
P1F-F does \textbf{not} support an inherent robustness advantage of the frozen literal AACOPF. Under the deliberately misspecified Gaussian likelihood, AACOPF is systematically less robust than conventional PF in the controlled campaign: it loses more clean successes, produces more wrong-mode locks, and generally has larger position and yaw errors.

The consequence for the thesis architecture is explicit. Robust UWB handling, if required in later phases, must be introduced as a separate mechanism---for example through a robust likelihood, innovation gating, an NLOS/reliability model, or reliability-aware measurement selection. The audited AACOPF copying transition must not itself be described as providing NLOS/outlier robustness.

No retuning of $(\alpha,\beta,c_\lambda)$ is performed in response to P1F-F. The next experiment, P1F-G, instead addresses the two remaining structural weaknesses already identified before this stress test: the quadratic all-pairs computational cost and the loss of rare useful support through aggressive particle cloning. Any future robust measurement model is kept conceptually separate from that scalable particle-transition adaptation.

The moderate/severe mixture cells are not strictly monotone in every empirical metric; for example, balanced PF correct-lock fraction is $0.90$ for the moderate case and $0.95$ for the severe case. With 20 seeds this cell-level ordering remains stochastic and is not interpreted as a benefit of larger outliers. The robust conclusion is based on the paired algorithm comparison across all stress cells.

The reproducible campaign is GitHub Actions workflow run \texttt{35000315361}. The aggregate artifact is \texttt{p1f-f-results}, artifact ID \texttt{10409367838}. Compact aggregate evidence is version controlled in \texttt{results/phase1/p1f\_f/}; full per-run evidence remains in the workflow artifact.
